% !TEX root = ../../../main.tex

\section{黑盒功能测试}

黑盒功能测试主要从系统使用者角度出发，
不关注程序内部实现细节，
而是通过页面操作、接口调用和业务流程执行结果，
验证系统功能是否符合需求分析和系统设计中的预期。

HearBridge 系统包含 HarmonyOS 移动端、Vue 管理端、Spring Boot 后端和 Python 手势识别服务等多个部分。

其中，
移动端主要面向普通用户，
用于完成登录态恢复、首页浏览、训练入口选择、手势资源列表查看、数字人动作展示、相机实时识别、raw 样本采集和句子视频识别等操作。

管理端主要面向管理员，
用于完成手势分类、手势资源、样本数据、训练任务和模型版本等资源维护工作。

Python 识别服务作为独立识别模块，
用于支撑移动端实时识别、raw dataset 采集、样本扫描、特征转换、模型训练、模型重载和句子视频识别等流程。

因此，
本节黑盒测试分别从移动端功能、后台管理端功能和跨端联动流程三个方面展开。

\subsection{移动端功能测试}

移动端功能测试主要验证普通用户在 HarmonyOS 应用中的核心使用流程。

根据当前移动端工程结构，
系统入口页面为 \texttt{Index}。

该页面在应用启动时恢复登录状态和最近练习记录，
并加载手语分类数据。

未登录时，
系统展示登录与注册相关页面；
登录后，
系统展示“首页”“训练”“我的”三个底部 Tab。

手语训练流程主要由 \texttt{GestureListPage} 和 \texttt{GesturePracticePageRefactored} 承担。

其中，
\texttt{GestureListPage} 根据分类编码加载手势资源列表，
支持分页加载更多资源，
并在点击资源后跳转至训练页面。

\texttt{GesturePracticePageRefactored} 负责手势训练页展示，
包含数字人播放、全屏相机 overlay、WebSocket 实时识别、当前模型版本展示以及开发者 raw dataset 采集模式。

句子视频识别流程由 \texttt{SignVideoRecognitionPage} 承担，
该页面支持选择本地视频、快速识别、词段 Top-K 详情展示以及 AI 忠实自然化结果展示。

表~\ref{tab:chapter07-mobile-black-box-test} 给出了移动端主要功能测试结果。

\begin{longtable}{L{0.08\textwidth} L{0.18\textwidth} L{0.24\textwidth} L{0.28\textwidth} L{0.12\textwidth}}
  \caption{移动端黑盒功能测试表}
  \label{tab:chapter07-mobile-black-box-test}\\
  \toprule
  编号 & 测试模块 & 测试内容 & 预期结果 & 测试结果 \\
  \midrule
  M01 & 应用启动 & 打开应用并进入 \texttt{Index} 页面 & 系统能够正常进入入口页面，并执行登录态恢复、最近练习恢复和分类加载逻辑 & 通过 \\
  M02 & 未登录展示 & 清除登录状态后启动应用 & 系统展示登录或注册相关页面，不直接进入主功能页面 & 通过 \\
  M03 & 登录态恢复 & 已登录状态下重新启动应用 & 系统恢复登录状态，并显示首页、训练、我的三个底部 Tab & 通过 \\
  M04 & 底部导航 & 点击首页、训练、我的三个 Tab & 页面能够正确切换，底部选中状态与展示内容保持一致 & 通过 \\
  M05 & 分类加载 & 启动后加载手语分类数据 & 首页和训练页能够展示分类数据；加载失败时显示错误提示 & 通过 \\
  M06 & 用户登录 & 输入正确用户信息后执行登录 & 登录成功后进入主页面，并保存用户登录状态 & 通过 \\
  M07 & 登录异常 & 输入错误账号、密码或验证码后执行登录 & 系统拒绝登录，并显示错误提示信息 & 通过 \\
  M08 & 我的页展示 & 进入我的页面查看用户信息 & 页面能够展示用户昵称、用户名、头像、分类数量和最近练习信息 & 通过 \\
  M09 & 最近练习 & 在我的页点击最近练习入口 & 若存在最近练习，系统跳转至对应手语训练页面；若不存在则不执行跳转 & 通过 \\
  M10 & 采集模式开关 & 在我的页开启或关闭采集模式 & 采集模式状态能够保存，并影响训练页中 raw dataset 采集入口与行为 & 通过 \\
  M11 & 手势资源列表 & 从训练页进入某一分类下的资源列表 & 页面展示该分类下资源名称、编码、封面和练习入口 & 通过 \\
  M12 & 资源分页加载 & 在资源列表中点击加载更多资源 & 系统请求下一页数据，并将新资源追加到当前列表中 & 通过 \\
  M13 & 进入训练页 & 点击某一手势资源的练习入口 & 系统跳转至 \texttt{GesturePracticePageRefactored}，并传入资源 ID、编码、中文名、封面和 SiGML 地址 & 通过 \\
  M14 & 数字人播放 & 在训练页点击播放或重新播放数字人动作 & 数字人区域能够加载 SiGML 资源并播放动作示范；异常时显示错误提示 & 通过 \\
  M15 & 相机 overlay & 在训练页点击打开相机 & 系统进入全屏相机 overlay，启动相机预览，并正确处理返回键关闭 overlay & 通过 \\
  M16 & 实时识别 & 在相机 overlay 中点击开始识别 & 系统通过 WebSocket 发送图像帧，并展示识别状态、识别结果、置信提示和当前模型版本 & 通过 \\
  M17 & 模型版本展示 & 实时识别会话建立后查看模型版本信息 & 页面显示当前实时识别服务使用的模型版本名称，并标识是否为管理端发布模型 & 通过 \\
  M18 & raw 样本采集 & 开启采集模式后在训练页采集当前动作样本 & 页面显示采集状态、有效帧数和采集窗口大小，并支持确认保存或丢弃待保存样本 & 通过 \\
  M19 & 视频选择 & 进入句子视频识别页并选择本地视频 & 系统打开视频选择器，选择成功后显示视频文件名，并重置历史识别结果 & 通过 \\
  M20 & 快速识别 & 选择视频后点击快速识别 & 系统上传视频并返回原始识别结果，包括原始中文、Gloss 序列、状态和耗时 & 通过 \\
  M21 & 词段详情 & 快速识别成功后查看词段详情 & 页面展示每个词段的标签、中文名、置信度、帧区间和 Top-K 候选 & 通过 \\
  M22 & AI 忠实自然化 & 快速识别后点击 AI 忠实自然化 & 页面展示候选排序结果、修正后 Gloss、中英文自然表达、置信等级和说明信息 & 通过 \\
  M23 & 句子识别异常 & 未选择视频时点击快速识别 & 系统提示“请先选择视频”，不发起无效识别请求 & 通过 \\
  M24 & 识别失败提示 & 视频识别或 AI 自然化接口异常 & 页面显示识别失败或自然化失败提示，不发生页面崩溃 & 通过 \\
  \bottomrule
\end{longtable}

由表~\ref{tab:chapter07-mobile-black-box-test} 可知，
移动端能够完成登录态恢复、底部 Tab 切换、手语分类加载、资源列表浏览、数字人动作展示、相机实时识别、raw 样本采集和句子视频识别等核心流程。

其中，
训练页验证了移动端与 Python 实时识别服务之间的 WebSocket 联动能力，
能够在相机预览过程中发送图像帧并展示识别状态与当前模型版本。

句子视频识别页验证了移动端短句视频识别与语义修正流程，
能够完成视频选择、快速识别、词段 Top-K 展示以及 AI 忠实自然化结果展示。

测试结果表明，
HarmonyOS 移动端已经能够支撑 HearBridge 系统普通用户侧的主要训练与识别功能。

\subsection{后台管理端功能测试}

后台管理端功能测试主要验证管理员是否能够通过 Web 页面完成资源维护和模型管理操作。

根据当前管理端工程结构，
系统通过 Vue Router 配置了后台登录页、Dashboard 页面和 HearBridge 管理菜单。

HearBridge 管理菜单下包含手势分类管理、手势资源管理、样本管理、训练管理和模型版本管理等页面。

路由守卫会在访问后台页面时检查本地管理员 token。

若用户未登录，
系统跳转至登录页；
若用户已登录，
则允许访问后台业务页面。

表~\ref{tab:chapter07-admin-black-box-test} 给出了后台管理端主要功能测试结果。

\begin{longtable}{L{0.08\textwidth} L{0.18\textwidth} L{0.24\textwidth} L{0.28\textwidth} L{0.12\textwidth}}
  \caption{后台管理端黑盒功能测试表}
  \label{tab:chapter07-admin-black-box-test}\\
  \toprule
  编号 & 测试模块 & 测试内容 & 预期结果 & 测试结果 \\
  \midrule
  A01 & 管理端登录 & 访问 \texttt{/login} 并输入正确管理员账号密码 & 登录成功后进入后台管理页面，并保存管理员 token & 通过 \\
  A02 & 登录异常 & 输入错误账号或密码后登录 & 系统拒绝登录，并显示登录失败提示 & 通过 \\
  A03 & 路由守卫 & 未登录状态下直接访问后台业务页面 & 系统跳转至登录页，并保留 redirect 参数 & 通过 \\
  A04 & 登录后跳转 & 已登录状态下访问登录页 & 系统自动跳转至 HearBridge 管理页面 & 通过 \\
  A05 & 侧边栏菜单 & 登录后查看 HearBridge 管理菜单 & 侧边栏展示手势分类、手势资源、样本、训练和模型版本等菜单项 & 通过 \\
  A06 & 手势分类列表 & 打开手势分类管理页面 & 页面能够展示分类 ID、分类编码、中文名称、封面和封面对象 Key & 通过 \\
  A07 & 手势分类新增 & 点击新增分类并填写分类编码、中文名称和封面对象 Key & 表单校验通过后保存分类，列表刷新并显示新增数据 & 通过 \\
  A08 & 手势分类编辑 & 点击分类行中的编辑按钮并修改名称或封面 & 系统保存修改结果，列表刷新后显示更新后的分类信息 & 通过 \\
  A09 & 手势分类删除 & 点击分类行中的删除按钮并确认 & 系统删除分类记录；若分类被资源引用，后端返回错误并由页面提示 & 通过 \\
  A10 & 分类分页 & 切换分类列表页码或每页数量 & 系统重新请求分页数据，并正确更新表格内容和分页状态 & 通过 \\
  A11 & 手势资源查询 & 在手势资源管理页面按分类筛选资源 & 表格返回指定分类下的手势资源，并展示编码、中文名、分类编码、封面和 SiGML 链接 & 通过 \\
  A12 & 手势资源新增 & 新增资源并填写编码、中文名、分类、封面和 SiGML 文件 & 系统保存资源信息，上传文件后返回对象 Key 和访问 URL & 通过 \\
  A13 & SiGML 上传校验 & 上传非 \texttt{.sigml} 或 \texttt{.xml} 文件，或上传超过限制的文件 & 页面拒绝上传，并显示格式或大小错误提示 & 通过 \\
  A14 & 手势资源编辑 & 编辑资源中文名、分类、封面或 SiGML 文件 & 系统保存修改结果，列表刷新后显示更新数据 & 通过 \\
  A15 & 手势资源删除 & 点击资源行中的删除按钮并确认 & 系统删除资源记录，列表刷新后该资源不再显示 & 通过 \\
  A16 & 样本统计展示 & 打开样本管理页面 & 页面展示总样本数、覆盖动作数、高质量样本、警告样本、异常样本和未知样本统计 & 通过 \\
  A17 & 样本筛选 & 按资源编码、质量状态或删除状态筛选样本 & 表格返回符合条件的样本记录，并展示帧数、时长、FPS、关键点比例和质量状态 & 通过 \\
  A18 & 同步样本 & 点击同步样本按钮并确认 & 系统从 Python 服务同步 raw dataset 样本，并刷新样本列表和统计数据 & 通过 \\
  A19 & 转换特征数据 & 点击转换特征数据按钮并确认 & 系统触发 raw dataset 到 feature 数据转换，并显示转换结果提示 & 通过 \\
  A20 & 样本质量标记 & 点击标记质量并填写质量状态和说明 & 系统保存样本质量状态，列表与统计信息同步刷新 & 通过 \\
  A21 & 样本删除 & 点击样本删除按钮并确认 & 系统执行样本软删除，列表刷新后样本删除状态更新 & 通过 \\
  A22 & 模型训练 & 打开训练管理页面并点击开始训练 & 系统触发模型训练任务，训练结束后展示训练名称、耗时、样本数、类别数、输入形状、准确率和产物路径 & 通过 \\
  A23 & 标签映射展示 & 模型训练结束后查看标签映射表格 & 页面展示每个手势标签对应的类别 ID & 通过 \\
  A24 & 模型版本列表 & 打开模型版本管理页面 & 页面展示版本名称、训练运行、样本数、类别数、验证准确率、训练耗时、状态和发布时间等信息 & 通过 \\
  A25 & 模型版本详情 & 点击模型版本详情按钮 & 弹窗展示基础信息、训练曲线、混淆矩阵、标签映射和评估结果等内容 & 通过 \\
  A26 & 模型版本发布 & 点击未发布版本的发布按钮并确认 & 系统将该版本设为当前发布版本，并刷新当前发布版本提示 & 通过 \\
  A27 & 当前发布版本 & 打开模型版本页面查看顶部发布状态 & 若存在发布版本，页面显示版本名称和验证准确率；否则显示尚未发布提示 & 通过 \\
  \bottomrule
\end{longtable}

由表~\ref{tab:chapter07-admin-black-box-test} 可知，
后台管理端能够支撑系统资源维护和模型管理流程。

手势分类管理页面能够完成分类新增、编辑、删除和分页展示；
手势资源管理页面能够完成按分类查询、资源新增编辑删除、封面上传和 SiGML 文件上传；
样本管理页面能够完成样本统计、条件筛选、同步样本、转换特征、质量标记和软删除；
训练管理页面能够触发模型训练并展示训练结果；
模型版本管理页面能够查看版本列表、查看训练产物详情并发布当前可用模型版本。

这些测试结果说明，
后台管理端已经能够支撑 HearBridge 手势识别数据管理和模型发布闭环。

\subsection{跨端联动流程测试}

跨端联动流程测试主要验证移动端、管理端、Spring Boot 后端和 Python 识别服务之间是否能够形成完整业务闭环。

与单一页面功能测试不同，
跨端联动流程需要多个模块共同完成。

例如，
手势资源需要先在管理端维护，
再由移动端加载使用；
raw 样本需要由移动端采集，
再由管理端同步、转换和训练；
模型版本需要在管理端发布，
再由 Python 服务重载，
最后由移动端实时识别页面使用。

表~\ref{tab:chapter07-integration-black-box-test} 给出了跨端联动流程测试结果。

\begin{longtable}{L{0.08\textwidth} L{0.20\textwidth} L{0.25\textwidth} L{0.29\textwidth} L{0.10\textwidth}}
  \caption{跨端联动黑盒功能测试表}
  \label{tab:chapter07-integration-black-box-test}\\
  \toprule
  编号 & 测试流程 & 测试内容 & 预期结果 & 测试结果 \\
  \midrule
  I01 & 分类资源联动 & 管理端新增手势分类和资源后，在移动端训练页查看分类与资源列表 & 移动端能够加载新增分类和资源，并可进入对应训练页面 & 通过 \\
  I02 & 资源文件联动 & 管理端上传资源封面和 SiGML 文件后，在移动端查看资源与数字人动作 & 移动端能够显示资源封面，并使用 SiGML 地址播放数字人动作 & 通过 \\
  I03 & 最近练习联动 & 移动端进入某一资源训练页后返回我的页 & 我的页能够显示最近练习信息，并支持再次进入该资源训练页 & 通过 \\
  I04 & raw 样本采集联动 & 移动端开启采集模式并采集当前动作样本 & Python 服务保存 raw dataset 样本，管理端后续可以同步到样本列表 & 通过 \\
  I05 & 样本同步联动 & 管理端点击同步样本按钮 & 后端调用 Python 服务扫描 raw 样本，样本管理页面刷新后显示新增或更新结果 & 通过 \\
  I06 & 特征转换联动 & 管理端点击转换特征数据按钮 & 后端调用 Python 服务执行 raw 到 feature 转换，并返回扫描、成功、跳过和失败数量 & 通过 \\
  I07 & 模型训练联动 & 管理端点击开始训练按钮 & Python 服务执行训练流程，后端记录训练结果，管理端展示模型路径、标签映射、训练曲线和混淆矩阵等信息 & 通过 \\
  I08 & 模型发布联动 & 管理端发布某一模型版本 & 后端更新发布状态，Python 服务可重载当前发布模型，移动端实时识别页能够显示对应模型版本 & 通过 \\
  I09 & 实时识别联动 & 移动端训练页打开相机并开始识别 & 移动端通过 WebSocket 向 Python 服务发送图像帧，并展示识别标签、置信度和模型版本 & 通过 \\
  I10 & 句子视频识别联动 & 移动端上传短句视频并点击快速识别 & Python 服务返回原始 Gloss 序列、中文映射和词段 Top-K，移动端展示原始识别结果和词段详情 & 通过 \\
  I11 & 语义修正联动 & 快速识别完成后点击 AI 忠实自然化 & 后端基于候选序列调用语义修正模块，移动端展示修正后 Gloss、中英文表达、候选编号和置信说明 & 通过 \\
  I12 & 异常联动处理 & 在识别服务异常、视频为空或请求参数错误时提交请求 & 系统返回错误提示，移动端或管理端页面不发生崩溃，并允许用户重新操作 & 通过 \\
  \bottomrule
\end{longtable}

由表~\ref{tab:chapter07-integration-black-box-test} 可知，
HearBridge 系统已经形成较完整的多端协同链路。

管理端维护的手势分类、手势资源和模型版本能够被移动端使用；
移动端采集的 raw 样本能够进入管理端样本管理和模型训练流程；
Python 识别服务能够支撑实时识别、样本采集、特征转换、模型训练、模型重载和句子视频识别；
Spring Boot 后端则在各端之间承担认证、数据管理、文件管理、模型版本记录和服务调用协调作用。

该测试结果说明，
系统并不是单独完成某个页面或某个算法实验，
而是已经能够完成从资源维护、样本采集、模型训练、模型发布到移动端识别展示的完整业务闭环。

\subsection{黑盒功能测试结果分析}

综合上述测试结果可知，
HearBridge 系统的移动端、后台管理端和跨端联动流程均能够完成主要功能测试。

移动端能够支撑普通用户完成登录态恢复、分类浏览、手势资源查看、数字人示范、实时识别、raw 样本采集和句子视频识别等操作。

后台管理端能够支撑管理员完成手势分类维护、手势资源维护、样本统计与筛选、样本同步、特征转换、模型训练和模型版本发布等操作。

跨端联动测试则验证了移动端、后端、管理端和 Python 识别服务之间的数据流转关系，
说明系统已经能够完成从样本采集到模型发布再到识别展示的闭环。

对于少数异常场景，
如无效登录、未授权访问、缺失参数、错误文件上传、未选择视频和识别服务异常等，
系统能够返回错误提示或拒绝访问结果，
没有出现严重崩溃问题。

因此，
黑盒功能测试结果表明，
HearBridge 系统已经具备毕业设计阶段所需的基本可用性和业务完整性。